% LaTeX source only. No TeX engine was run when producing this document.
\documentclass[UTF8,a4paper,11pt,fontset=fandol]{ctexart}
\usepackage[margin=2.2cm]{geometry}
\usepackage{amsmath,amssymb,bm,booktabs,graphicx}
\usepackage{xcolor,listings,hyperref,url}
\hypersetup{colorlinks=true,linkcolor=blue,urlcolor=blue}
\lstset{language=Python,basicstyle=\ttfamily\small,breaklines=true,
  frame=single,columns=fullflexible,showstringspaces=false}
\newcommand{\CI}{\mathrm{CI}}
\newcommand{\T}{\mathsf{T}}
\newcommand{\vrow}{\operatorname{vec}_{\mathrm C}}
\title{FIC--NEVPT2 解析梯度中的 CI 源项收缩优化\\推导、实现与前后性能对照}
\author{基于本地代码和实测数据的学习笔记}
\date{2026 年 10 月 18 日}
\begin{document}
\maketitle
\begin{abstract}
本笔记说明如何将原来逐 CI 系数生成过渡高阶密度矩阵的源项计算，
改写为压缩多电子湮灭算符对整个 CI 向量的作用。
旧算法完整保留为 \texttt{legacy} 后端，新算法为 \texttt{contracted} 后端。
这是一项精确的代数重组，不改变 FIC--NEVPT2 的能量定义、CI 空间或收缩近似。
本文件只提供 LaTeX 源代码，本次未编译。
\end{abstract}
\tableofcontents

\section{优化的是梯度的哪一部分}
设参考波函数为
\[
 |\Psi\rangle=\sum_{I=1}^{D}c_I|I\rangle,\qquad c^\T c=1.
\]
NEVPT2 梯度代码已经将微扰 Lagrangian 对 RDM 的导数累积为权重
\[
 W^{(k)}=\frac{\partial\mathcal L}{\partial\Gamma^{(k)}},
 \qquad k=1,2,3,4.
\]
在源代码中，它们分别保存在
\begin{lstlisting}
self.rdms.ca
self.rdms.ccaa
self.rdms.cccaaa
self.rdms.ccccaaaa
\end{lstlisting}
需要计算的 RDM 部分 CI 源项为
\begin{equation}
 b_I^{\CI,\mathrm{RDM}}
 =\sum_{k=1}^4 W^{(k)}:
       \frac{\partial\Gamma^{(k)}}{\partial c_I}.
 \label{eq:chain}
\end{equation}
这里 $W^{(k)}$ 是当前参考点上已经求出的偏导权重；
应用链式法则时不应再对这些权重重复求导。
这并不意味着完整 Lagrangian 就等于 $\sum_k W^{(k)}:\Gamma^{(k)}$；
该表达式只是用于获得式 \eqref{eq:chain} 的辅助收缩。

本次只替换式 \eqref{eq:chain} 的计算。Fock/轨道能量带来的额外 CI 源项、
其他态的贡献、半正则轨道响应、CI 归一化投影和 Z-vector 求解仍沿用原流程。

\section{旧实现：逐 CI 单位向量求过渡 RDM}
对每个 $I$ 构造 $e_I$，再计算 $\langle I|\hat O^{(k)}|\Psi\rangle$。
实数波函数下还需加上 bra/ket 交换的贡献：
\[
 \frac{\partial\Gamma^{(k)}_{PQ}}{\partial c_I}
 =\langle I|\hat O^{(k)}_{PQ}|\Psi\rangle
  +\langle\Psi|\hat O^{(k)}_{PQ}|I\rangle.
\]
原实现的基本结构为
\begin{lstlisting}
for I in range(state_ci.size):
    unit_ci.flat[I] = 1.0
    transition_rdms = make_tran_rdm1234(unit_ci, state_ci, ...)
    # Add the bra/ket transpose and contract with RDM weights.
    source.flat[I] += weighted_transition_contractions
    unit_ci.flat[I] = 0.0
\end{lstlisting}
因此会进行 $D$ 次过渡 1--4-RDM 构造。总时间可以写成
\[
 T_{\rm old}\sim
 D\,T_{\rm transition\ RDM}(a,D)+D\,T_{\rm RDM\ contraction}(a),
\]
其中 $a$ 为活性轨道数。
旧算法不一定同时保存 $D$ 个 4-RDM；主要问题是重复生成、补全、转置和收缩。

\section{从正常序 RDM 到湮灭后的 CI 向量}
令 $P=(p_1,\ldots,p_k)$、$Q=(q_1,\ldots,q_k)$，自旋模式为
$\sigma=(\sigma_1,\ldots,\sigma_k)$。正常序、自旋求和的 RDM 为
\[
 \Gamma^{(k)}_{PQ}
 =\sum_\sigma
 \langle\Psi|
 a^\dagger_{p_1\sigma_1}\cdots a^\dagger_{p_k\sigma_k}
 a_{q_k\sigma_k}\cdots a_{q_1\sigma_1}
 |\Psi\rangle.
\]
定义湮灭后的向量
\[
 |T_{Q\sigma}\rangle
 =a_{q_k\sigma_k}\cdots a_{q_1\sigma_1}|\Psi\rangle.
\]
对剩余电子的行列式 $|J\rangle$ 展开，有
\[
 \Gamma^{(k)}_{PQ}
 =\sum_{\sigma,J}T_{P\sigma,J}T_{Q\sigma,J}.
\]
于是 RDM 可以视为湮灭后 CI 向量的 Gram 收缩，
无需为每个原始 CI 单位向量单独计算过渡 RDM。

\section{自旋分块和反对称轨道组合压缩}
对 $k=m_\alpha+m_\beta$，把相同自旋计数的模式归为一块。
每块只保存递增的同自旋轨道组合，维数为
\[
 M_s=\binom{a}{m_\alpha}\binom{a}{m_\beta},\qquad
 R_s=\binom{a}{n_\alpha-m_\alpha}
     \binom{a}{n_\beta-m_\beta}.
\]
$M_s$ 是压缩轨道元组数，$R_s$ 是湮灭后 CI 空间的维数。
超出实际电子数的块直接跳过。

设 $\widetilde W_s$ 为所有相应自旋排列和费米反对称符号合并后的权重。
\textbf{使用组合选基不等于丢弃排列贡献}：
代码显式把所有排列按正负号累积到 $\widetilde W_s$，
因此不应再额外手动乘除 $k!$。

例如 $a=8,k=4,m_\alpha=m_\beta=2$ 时，
\[
 a^4=4096,\qquad M_s=\binom82^2=784.
\]
完整四体权重仍有 $a^8$ 元素，但该自旋块的作用可在压缩空间中进行。

\subsection{稀疏湮灭算符}
构造
\[
 A_s\in\mathbb R^{(M_sR_s)\times D},\qquad
 T_s=\operatorname{reshape}(A_s c,M_s,R_s).
\]
$A_s$ 只包含允许的占据数连接和 $\pm1$ 的费米符号，采用 CSR 存储，
不构造这个矩形矩阵的稠密表示。
本项目 CI 下标按 alpha 字符串优先、beta 字符串次之排列：
\[
 I=I_\alpha N_\beta+I_\beta.
\]
符号由 \texttt{cistring.des\_sign} 求出。同一压缩自旋块内，
固定的 alpha/beta 交叉相位会在 $A_s$ 与 $A_s^\T$ 之间抵消。

全 CI 空间中，每块的非零连接数为
\[
 \operatorname{nnz}(A_s)
 =D\binom{n_\alpha}{m_\alpha}\binom{n_\beta}{m_\beta}.
\]
这里仍需要访问 CI 系数和行列式连接；优化消除的是 $D$ 次完整过渡 RDM，
不是消除 CI 空间本身。

\section{直接计算整个 CI 源项}
固定权重时，每块辅助收缩为
\[
 \mathcal F_s
 =\sum_{PQJ}\widetilde W_{s,PQ}T_{s,PJ}T_{s,QJ}.
\]
求导得到
\begin{align}
 Y_s&=(\widetilde W_s+\widetilde W_s^\T)T_s,\\
 \bm b^{\CI,\mathrm{RDM}}
 &=\sum_s A_s^\T\vrow(Y_s).
 \label{eq:action}
\end{align}
$\vrow$ 表示与 NumPy 的 C-order \texttt{ravel()} 一致的展开。
若写成单个线性算符，式 \eqref{eq:action} 等价于
\[
 A_s^\T\big[(\widetilde W_s+\widetilde W_s^\T)
                  \otimes I_{R_s}\big]A_s c,
\]
但代码不形成 Kronecker 积。

核心实现恰好是：
\begin{lstlisting}
op, residual = _annihilator(norb, nelec, na, nb, tuples)
ket = (op @ ci.ravel()).reshape(count, residual)
compressed = _compressed_weight(weight, norb, rank, na, nb, tuples)
acted = (compressed + compressed.T) @ ket
out += op.T @ acted.ravel()
\end{lstlisting}
bra 和 ket 两侧的导数由 $\widetilde W+\widetilde W^\T$ 一起处理。

\subsection{与 NEVPT2 代码的对应}
\begin{itemize}
 \item \path{grad/ci_contraction.py}：
 \texttt{\_des\_map} 生成连接与符号；
 \texttt{\_annihilator} 组装 CSR；
 \texttt{\_compressed\_weight} 合并自旋和反对称排列；
 \texttt{contract\_rdm\_weights} 执行式 \eqref{eq:action}。
 \item \path{grad/gradients.py}：
 \texttt{get\_ci\_response\_legacy()} 完整保留旧循环；
 \texttt{get\_ci\_response\_contracted()} 调用新收缩；
 \texttt{get\_ci\_response()} 根据后端分派并记录统计。
 \item \path{grad/dfgradients.py}：
 保留 DF 原有旧循环，复用相同的新源项接口；
 DF 积分和辅助基响应的其余部分未改变。
\end{itemize}
得到 CI 源项后，原程序继续追加其余源项并做归一化投影，
再与轨道部分组成耦合 Z-vector 方程。没有提前改变投影顺序。

\section{如何切换并保留旧实现}
新默认后端为 \texttt{contracted}，旧路径仍可显式调用。
对相同已收敛参考态，分别构造新的 PT/梯度对象：
\begin{lstlisting}
pt_old = NEVPT2(mc, ...)
pt_old.kernel()
g_old = pt_old.nuc_grad_method()
g_old.ci_response_backend = 'legacy'
gradient_old = g_old.kernel()

pt_new = NEVPT2(mc, ...)
pt_new.kernel()
g_new = pt_new.nuc_grad_method()
g_new.ci_response_backend = 'contracted'
gradient_new = g_new.kernel()
print(g_new.ci_response_stats)
\end{lstlisting}
不要在同一个已累积过中间量的梯度对象上重复运行完整 kernel 来计时。
常规积分和 DF 梯度都支持该开关。

\section{测试条件和结果一致性}
测试在 Intel Core i7-10700、约 16 GB 内存的本机执行，
OpenMP/OpenBLAS/MKL 均固定为一个线程。
每组新旧后端共享同一已收敛参考态，使用新的 PT 和梯度对象，
先计算能量再计时梯度；三次重复并交替执行顺序，报告中位数。
辅助模块在计时前导入；没有把剖析开销计入普通时间。

常规积分参考使用 Newton-CASSCF；DF 参考使用标准 DF-CASSCF，
辅助基为 \texttt{weigend}，包含辅助基响应。
单态参考收敛后包装为 NEVPT2 所需的单根容器，没有额外优化其他态；
双态平均使用等权重。测试工具包含当前 PySCF 版本的单根 list/array 形状适配，
与本次源项代数修改分开。

\begin{table}[htbp]\centering\small
\caption{算例。CAS 写成（电子数，轨道数）；目标根从 0 开始。}
\begin{tabular}{lllrrrrr}\toprule
编号 & 分子/基组 & CAS & $D$ & 积分 & 根数 & 目标根 & canonstep\\\midrule
C1 & LiH/STO-3G & $(2,2)$ & 4 & ERI & 1 & 0 & 0 \\
C2 & H$_2$O/STO-3G & $(2,2)$ & 4 & ERI & 1 & 0 & 0 \\
C3 & LiH/STO-3G & $(4,4)$ & 36 & ERI & 1 & 0 & 0 \\
C4 & LiH/STO-3G & $(4,5)$ & 100 & ERI & 1 & 0 & 0 \\
C5 & LiH/6-31G & $(4,6)$ & 225 & ERI & 1 & 0 & 0 \\
C6 & H$_2$O/STO-3G & $(2,2)$ & 4 & DF & 1 & 0 & 0 \\
C7 & LiH/STO-3G & $(4,5)$ & 100 & DF & 1 & 0 & 0 \\
C8 & LiH/STO-3G & $(2,2)$ & 4 & ERI & 2 & 0 & 2 \\
C9 & LiH/STO-3G & $(2,2)$ & 4 & ERI & 2 & 1 & 2 \\
C10 & LiH/STO-3G & $(2,2)$ & 4 & ERI & 1 & 0 & 1 \\
\bottomrule\end{tabular}
\end{table}

对全部十组测试，前后最大差异为：
\begin{align*}
 |\Delta E|_{\max}&=0,\\
 \|\Delta b^{\CI,\mathrm{RDM}}\|_\infty
 &=5.004\times10^{-15},\\
 \|\Delta b_{\rm full}\|_\infty
 &=5.004\times10^{-15},\\
 \|\Delta Z\|_\infty&=1.306\times10^{-14},\\
 \|\Delta\nabla E\|_\infty&=6.211\times10^{-16}
 \quad\text{Hartree/Bohr}.
\end{align*}
核梯度前后比较覆盖每个算例的全部坐标分量。
此外每个算例对末原子 $z$ 坐标做全电子总能中心差分抽查：
\[
 g_{\rm num}=\frac{E(R+h)-E(R-h)}{2h},\qquad h=5\times10^{-4}\ {\rm Bohr}.
\]
每个位移点重新收敛参考态和 NEVPT2，包含核排斥、轨道和 CI 的变化。
两种后端的最大抽查误差均不超过
$9.031\times10^{-8}$ Hartree/Bohr。
这是代表坐标的差分抽查，不应称为所有坐标的数值差分验证。

\subsection{前后耗时}
\begin{table}[htbp]\centering\small
\caption{三次计时中位数，时间单位为秒。CI 源项时间和完整梯度时间分别报告。}
\resizebox{\textwidth}{!}{%
\begin{tabular}{lrrrrrr}\toprule
编号 & CI 旧 & CI 新 & CI 加速 & 梯度旧 & 梯度新 & 梯度加速\\\midrule
C1 & 0.004739 & 0.001086 & 4.4 & 0.1759 & 0.1605 & 1.10 \\
C2 & 0.004839 & 0.001097 & 4.4 & 0.1790 & 0.1768 & 1.01 \\
C3 & 0.266768 & 0.005605 & 47.6 & 0.3571 & 0.0969 & 3.68 \\
C4 & 1.805743 & 0.011789 & 153.2 & 2.0010 & 0.1915 & 10.45 \\
C5 & 11.347608 & 0.032685 & 347.2 & 11.8778 & 0.5505 & 21.58 \\
C6 & 0.004679 & 0.001079 & 4.3 & 0.2303 & 0.2148 & 1.07 \\
C7 & 1.817032 & 0.011385 & 159.6 & 2.0382 & 0.2315 & 8.80 \\
C8 & 0.004765 & 0.001091 & 4.4 & 0.1690 & 0.1675 & 1.01 \\
C9 & 0.004693 & 0.001070 & 4.4 & 0.1715 & 0.1672 & 1.03 \\
C10 & 0.004733 & 0.001104 & 4.3 & 0.1781 & 0.1621 & 1.10 \\
\bottomrule\end{tabular}}
\end{table}

对 C5（LiH/CAS(4,6)，225 CI 分量），CI 源项由约
11.348 秒降到 0.0327 秒，加速约 347 倍；
完整梯度由 11.878 秒降到 0.5505 秒，加速约 21.6 倍。
旧路径对该源项构造 225 次过渡高阶 RDM，新路径为零次。
仍需一次正常前向 RDM；Z-vector 内部也仍有低阶过渡密度计算。

\subsection{为什么完整梯度加速小于源项加速}
若源项占原总时间的比例为 $f$，源项加速为 $S_{\CI}$，则近似有
\[
 S_{\rm total}\approx
 \frac{1}{(1-f)+f/S_{\CI}}.
\]
C5 中 $f\approx0.955$，该估计给出约 21 倍，与实测约 21.6 倍相符。
小 CAS 中该源项只占几个百分点，因此即使源项加速数倍，
完整梯度也可能只略有变化。缓存、库初始化和小样本计时波动也会影响小算例。

\section{完整结构优化前后对照}
同样采用 geomeTRIC、相同初始构型、相同参考态更新及收敛阈值，
两种后端分别真正运行到收敛：
\begin{table}[htbp]\centering\small
\caption{优化总墙钟时间包含参考态更新、PT 能量、梯度和优化器开销。}
\begin{tabular}{lrrrr}\toprule
体系 & 梯度评估次数旧/新 & 旧时间/秒 & 新时间/秒 & 加速\\\midrule
H$_2$O/STO-3G, CAS(2,2) & 5/5 & 1.6089 & 1.5471 & 1.04 \\
LiH/STO-3G, CAS(4,5) & 4/4 & 8.5556 & 1.3072 & 6.54 \\
\bottomrule\end{tabular}
\end{table}
最终坐标最大分量差异分别约为 $5.05\times10^{-12}$ 和
$2.04\times10^{-11}$ Bohr；终点能量差异分别约为
$2.84\times10^{-14}$ 和 $2.66\times10^{-15}$ Hartree。
这些差异处在本测试的浮点误差范围内，不代表引入了新的电子结构近似。

\section{时间和内存标度仍有哪些限制}
新源项的主要成本可按自旋块写成
\[
 T_{\rm new}\sim
 \sum_s\left[
 O(\operatorname{nnz} A_s)+O(M_s^2R_s)
 \right]+T_{\rm compress}(a).
\]
这消除了逐 CI 单位向量的完整过渡 RDM 调用，但 CI 维数仍是
\[
 D=\binom{a}{n_\alpha}\binom{a}{n_\beta},
\]
半填充时仍有组合增长。
权重和正常前向 4-RDM 的 $a^8$ 存储、活性收缩块、
CASSCF 与 Z-vector 迭代、外部积分和 AO 有效密度仍可能成为瓶颈。

一个源项自旋块的存储约为
\[
 O(\operatorname{nnz}A_s+M_sR_s+M_s^2),
\]
并且一次只处理一个块。不能把旧算法的峰值内存误写成 $D$ 个 4-RDM
同时驻留，也不能把此次加速解读为消除了所有大 CAS 或大分子限制。

\section{复现、回归与适用边界}
\begin{enumerate}
 \item \path{tests/test_ci_response.py} 对常规/DF 两个梯度类，
 在随机 RDM 权重和不同 alpha/beta 电子数下比较新旧源项；
 禁止新路径调用过渡 4-RDM，检查内存预算和错误后端处理。
 这些不同自旋的单元测试不等于开壳层分子全部梯度路径验证。
 \item 成对基准保存源项、完整右端、Z-vector、核梯度、耗时和
 正负位移能量；源文件的归一化 LF SHA256 也记录在每份 JSON 中。
 \item 三次重复按新旧交替顺序串行执行，避免把不同机器、不同参考态
 或不相同的计算设置直接相除。峰值 RSS 为整个进程的峰值，
 不能据此单独归因每个后端的增量内存。
 \item 单根容器兼容适配、严格为空的激发类别筛选均位于基准工具中，
 没有混入本次 NEVPT2 源项修改。
 \item 验证包括常规积分、DF、canonstep=0/1/2、等权双态平均的两个目标根。
 没有据此声称验证了所有版本、复数波函数或所有开壳层分子。
\end{enumerate}

示例命令（在本机 PySCF 工作区运行，按实际安装环境调整路径）：
\begin{lstlisting}[language=bash]
export PYTHONPATH="$PWD"
export OMP_NUM_THREADS=1 OPENBLAS_NUM_THREADS=1 MKL_NUM_THREADS=1
python tools/benchmark_nevpt_ci_response.py \
  --nevpt-path ../nevpt_grad/embed_sim \
  --case lih_631g_e4a6 --repeats 3 --fd-check
python tools/run_nevpt_ci_response_benchmarks.py \
  --nevpt-path ../nevpt_grad/embed_sim
python tools/run_nevpt_ci_response_benchmarks.py \
  --nevpt-path ../nevpt_grad/embed_sim --optimization-only
\end{lstlisting}

\section{学习时建议核对的三件事}
第一，确认权重对应的是正常序、创建/湮灭指标分组后的 RDM；
NEVPT2 与通用 PySCF 的 1-RDM 返回约定可能相差转置，
这里使用实数波函数并同时处理 bra/ket 导数，必须通过旧实现逐项核验。
第二，明确组合压缩的排列符号已经吸收到权重，不能再添加不一致的阶乘。
第三，把原始 CI 源项、完整投影右端和最终核梯度分层比较；
仅比较最终优化结构接近，不足以证明导数实现正确。
\end{document}
